% Created 2024-01-08 lun 09:11
% Intended LaTeX compiler: pdflatex
\documentclass[11pt]{article}
\usepackage[utf8]{inputenc}
\usepackage[T1]{fontenc}
\usepackage{graphicx}
\usepackage{longtable}
\usepackage{wrapfig}
\usepackage{rotating}
\usepackage[normalem]{ulem}
\usepackage{amsmath}
\usepackage{amssymb}
\usepackage{capt-of}
\usepackage{hyperref}
\usepackage{minted}

\usepackage{parskip}
\usepackage[utf8]{inputenc} %% For unicode chars
\usepackage{placeins}
\usepackage[margin=2.50cm]{geometry}
\usepackage[T1]{fontenc}
\usepackage{mathpazo}
\linespread{1.05}
\usepackage[scaled]{helvet}
\usepackage{courier}
\hypersetup{colorlinks=true,linkcolor=blue}
\RequirePackage{fancyvrb}
\DefineVerbatimEnvironment{verbatim}{Verbatim}{fontsize=\small,formatcom = {\color[rgb]{0.5,0,0}}}
\usepackage{mdframed}
\BeforeBeginEnvironment{minted}{\begin{mdframed}}
\AfterEndEnvironment{minted}{\end{mdframed}}
\setcounter{secnumdepth}{3}
\date{}
\title{PHP Sessions and Cookies}
\hypersetup{
 pdfauthor={},
 pdftitle={PHP Sessions and Cookies},
 pdfkeywords={},
 pdfsubject={},
 pdfcreator={Emacs 28.1 (Org mode 9.5.5)}, 
 pdflang={Spanish}}
\begin{document}

\maketitle
\setcounter{tocdepth}{3}
\tableofcontents

\setlength\parindent{10pt}

\section{Sessions And Cookies}
\label{sec:orge4e42ef}
Managing sessions and cookies is an essential aspect of web development
in PHP, as they allow you to maintain state across multiple page
requests. Here's a guide to help you understand and use sessions and
cookies effectively in PHP.

\subsection{PHP Sessions}
\label{sec:org9f1dfaa}
Sessions in PHP are used to preserve certain data across subsequent
accesses by the same user. This is useful for things like maintaining
user state (e.g., ensuring a user is logged in) across different pages.

\subsubsection{Starting a Session}
\label{sec:orgf9c3620}
To start a session, use \texttt{session\_start()}. This function must be called
before any output is sent to the browser.

\begin{minted}[]{php}
session_start();
\end{minted}

\subsubsection{Storing Data in a Session}
\label{sec:org9b417ed}
Once a session is started, you can store data in the \texttt{\$\_SESSION}
superglobal array.

\begin{minted}[]{php}
$_SESSION["username"] = "john_doe";
\end{minted}

\subsubsection{Accessing Session Data}
\label{sec:org7af6b31}
You can access session data stored in the \texttt{\$\_SESSION} array as long as
the session is active.

\begin{minted}[]{php}
echo "Welcome, " . $_SESSION["username"];
\end{minted}

\subsubsection{Ending a Session}
\label{sec:org89c14b4}
To end a session and clear session data, use \texttt{session\_unset()} and
\texttt{session\_destroy()}.

\begin{minted}[]{php}
session_unset(); // remove all session variables
session_destroy(); // destroy the session
\end{minted}

\subsection{PHP Cookies}
\label{sec:orgf9a2e25}
Cookies are small files stored on the user's computer. They are
typically used to remember information about users, such as login
details or website preferences.

\subsubsection{Setting Cookies}
\label{sec:orgfe84655}
Use the \texttt{setcookie()} function to set a cookie. This function must be
called before any output is sent to the browser.

\begin{minted}[]{php}
setcookie("user", "john_doe", time() + 3600); // 3600 seconds = 1 hour
\end{minted}

\subsubsection{Accessing Cookies}
\label{sec:org70d8ce5}
Access cookies through the \texttt{\$\_COOKIE} superglobal array.

\begin{minted}[]{php}
if(isset($_COOKIE["user"])) {
    echo "Welcome " . $_COOKIE["user"];
}
\end{minted}

\subsubsection{Deleting Cookies}
\label{sec:orgf5898fe}
Delete a cookie by setting its expiration date in the past.

\begin{minted}[]{php}
setcookie("user", "", time() - 3600);
\end{minted}

\subsection{Best Practices and Security Considerations}
\label{sec:org52bc7e4}
\subsubsection{Sessions}
\label{sec:org97fc87b}
\begin{itemize}
\item Regenerate session ID: Use \texttt{session\_regenerate\_id()} to prevent
session fixation attacks.
\item Secure session cookies: Set \texttt{session.cookie\_secure} and
\texttt{session.cookie\_httponly} flags to true to enhance security.
\item Session data is stored server-side, which is generally more secure
than storing sensitive data in cookies.
\end{itemize}

\subsubsection{Cookies}
\label{sec:org08a27a2}
\begin{itemize}
\item Sensitive data: Avoid storing sensitive data in cookies as they are
stored client-side.
\item Secure and HTTPOnly flags: Set these flags when setting cookies to
enhance security (\texttt{secure} for HTTPS only, \texttt{HTTPOnly} to make them
inaccessible to JavaScript).
\item Expiration time: Be mindful of the expiration time of cookies,
especially for sensitive information.
\end{itemize}

\subsubsection{General}
\label{sec:orgae30cf7}
\begin{itemize}
\item Data validation: Always validate and sanitize any data received from
cookies or sessions.
\item HTTPS: Use HTTPS to encrypt cookies and session data transmitted over
the network.
\end{itemize}

Using sessions and cookies appropriately can greatly enhance the user
experience on your website by allowing for persistent, stateful
interactions across multiple web pages. Remember to handle them securely
to protect your users' data.

\section{Sessions and cookies cont.}
\label{sec:org8f2d84c}
Certainly! Below are code examples to illustrate how to use cookies and
sessions in PHP. These examples cover basic operations like setting,
accessing, and deleting cookies, as well as starting a session, storing
session data, and ending a session.

\subsubsection{Using Cookies in PHP}
\label{sec:org0825c7a}
\begin{itemize}
\item Setting a Cookie
\label{sec:orgae862f4}
To set a cookie, use the \texttt{setcookie()} function. This should be done
before any output is sent to the browser.

\begin{minted}[]{php}
// Set a cookie that expires in 1 hour
setcookie("username", "john_doe", time() + 3600, "/");

// Secure and HTTPOnly flags
setcookie("secure_user", "john_doe_secure", time() + 3600, "/", "", true, true);
\end{minted}

\item Accessing a Cookie
\label{sec:org0ff0dc2}
Access a cookie via the \texttt{\$\_COOKIE} superglobal array.

\begin{minted}[]{php}
if (isset($_COOKIE["username"])) {
    echo "Welcome " . htmlspecialchars($_COOKIE["username"]);
} else {
    echo "Welcome guest!";
}
\end{minted}

\item Deleting a Cookie
\label{sec:orgec7c807}
Delete a cookie by setting its expiration date to a past time.

\begin{minted}[]{php}
// Deleting a cookie by setting its expiration to one hour ago
setcookie("username", "", time() - 3600, "/");
\end{minted}
\end{itemize}

\subsubsection{Using Sessions in PHP}
\label{sec:org957c1cd}
\begin{itemize}
\item Starting a Session
\label{sec:org9362221}
Start a session with \texttt{session\_start()}. This should be the first thing
in your script.

\begin{minted}[]{php}
session_start();
\end{minted}

\item Storing Data in a Session
\label{sec:orgf8e16eb}
Store data in the \texttt{\$\_SESSION} superglobal array after starting the
session.

\begin{minted}[]{php}
$_SESSION["username"] = "john_doe";
\end{minted}

\item Accessing Session Data
\label{sec:org1b62024}
Retrieve data from the session.

\begin{minted}[]{php}
session_start();
if (isset($_SESSION["username"])) {
    echo "Welcome " . htmlspecialchars($_SESSION["username"]);
} else {
    echo "Welcome guest!";
}
\end{minted}

\item Ending a Session
\label{sec:org9793bb9}
End a session and clear its data.

\begin{minted}[]{php}
session_start();

// Unset all session variables
session_unset();

// Destroy the session
session_destroy();
\end{minted}
\end{itemize}

\subsubsection{Important Notes}
\label{sec:org415a6c0}
\begin{itemize}
\item \textbf{Cookie Security}: Remember that cookies are stored client-side. Avoid
storing sensitive information in them. Use the \texttt{secure} and \texttt{HTTPOnly}
flags for enhanced security, especially if using cookies for
authentication.

\item \textbf{Session Security}: PHP sessions are generally more secure as the
session data is stored server-side. However, ensure to regenerate
session IDs (\texttt{session\_regenerate\_id()}) during login to prevent
session fixation attacks.

\item \textbf{Output Buffering}: If you need to set cookies after sending output,
use output buffering (\texttt{ob\_start()}) at the beginning of your script.
\end{itemize}

These examples cover fundamental uses of cookies and sessions in PHP.
For more advanced features and security measures, consult the official
PHP documentation and consider additional security practices relevant to
your application's context.
\end{document}
